\chapter{Thrombotic Thrombocytopenic Purpura}
\index{Thrombotic Thrombocytopenic Purpura}
\label{sec:Thrombotic Thrombocytopenic Purpura}

\section{Overview}
Thrombotic thrombocytopenic purple spots on the skin (TTP) is a rare, life-threatening thrombotic microangiopathy characterized by microvascular blood clot formation, thrombocytopenia, and microangiopathic hemolytic anemia, with an incidence of 3--4 per million. This condition affects a significant number of people worldwide and can range from mild to severe in presentation. Prompt diagnosis and appropriate management are essential for optimal outcomes. Early recognition and appropriate management are essential for optimal outcomes. A thorough understanding of the underlying mechanisms guides treatment decisions. Patient education and self-management play important roles in long-term care.
\section{Causes}
This condition happens because of severe deficiency of ADAMTS13, a von Willebrand factor (VWF) cleaving protease: immune-mediated (acquired TTP, IgG autoantibodies, adulthood) or congenital (Upshaw-Schulman syndrome, ADAMTS13 gene mutations, childhood). ADAMTS13 deficiency leads to accumulation of ultralarge VWF multimers, causing spontaneous platelet adhesion and microthrombi formation in arterioles and capillaries. The underlying mechanisms involve complex interactions between genetic predisposition and environmental factors that disrupt normal physiological processes. The underlying mechanisms involve complex interactions between genetic predisposition and environmental factors. Disruption of normal physiological processes leads to the characteristic features of the condition. Multiple contributing factors may act synergistically to trigger or worsen the condition. Understanding the specific cause helps guide targeted treatment approaches.
\section{Risk Factors}
Genetic predisposition and family history Advancing age Lifestyle factors (diet, exercise, smoking, alcohol) Environmental exposures Underlying medical conditions Medication use Stress and psychological factors Socioeconomic factors

\begin{itemize}[noitemsep]
  \item Genetic predisposition and family history
  \item Advancing age
  \item Lifestyle factors (diet, exercise, smoking, alcohol)
  \item Environmental exposures
  \item Underlying medical conditions
  \item Medication use
\end{itemize}
\section{Signs and Symptoms}
\begin{itemize}[noitemsep]
  \item You may experience the classic pentad: thrombocytopenia (tiny red or purple spots, purple spots on the skin, bleeding), microangiopathic hemolytic anemia (fever, pallor, jaundice), nervous system involvement (headache, confusion, seizures, focal deficits, coma), kidney involvement (elevated creatinine, hematuria), and fever. Not all five features are present simultaneously
  \item Thrombocytopenia and MAHA are sufficient to suspect TTP.
  \item Pain or discomfort in the affected area
  \item Functional impairment affecting daily activities
  \item Changes in appearance or function of affected tissues
  \item Systemic symptoms may include fatigue, fever, or weight changes
  \item Symptoms may develop gradually or appear suddenly
\end{itemize}
\section{Progression and Stages}
The condition may progress through different stages of severity over time. Early stages often involve mild or subtle symptoms that may be overlooked. Moderate stages involve more noticeable symptoms affecting daily function. Advanced stages may involve significant impairment and complications.
\section{Complications}
\begin{itemize}[noitemsep]
  \item Disease progression leading to worsening symptoms
  \item Secondary infections or injuries
  \item Side effects of treatment
  \item Reduced quality of life
  \item Emotional and psychological impact
\end{itemize}
\section{Diagnosis}
Doctors diagnose this based on peripheral blood smear showing schistocytes (helmet cells, fragmented RBCs), elevated LDH, low haptoglobin, negative direct Coombs test, and ADAMTS13 activity <10\%. Comprehensive medical history and physical examination Laboratory tests to assess organ function and rule out other conditions Imaging studies to visualize affected structures Specialized testing depending on the affected system Consultation with relevant specialists Monitoring of disease progression over time

\begin{itemize}[noitemsep]
  \item Comprehensive medical history and physical examination
  \item Laboratory tests to assess organ function
  \item Imaging studies to visualize affected structures
  \item Specialized testing depending on the affected system
  \item Consultation with relevant specialists
\end{itemize}
\section{Prevention}
\begin{itemize}[noitemsep]
  \item Maintain a healthy lifestyle with balanced diet and exercise
  \item Avoid known risk factors
  \item Attend regular medical check-ups for early detection
  \item Follow recommended screening guidelines
\end{itemize}
\section{Treatment Options}
Treatment typically involves a medical emergency: urgent plasma exchange (PLEX, 1--1.5 plasma volume daily) is the cornerstone of therapy, combined with high-dose corticosteroids (prednisone 1 mg/kg). For refractory/relapsed immune TTP: rituximab, caplacizumab (anti-VWF nanobody, blocks platelet adhesion), or immunosuppression (vincristine, cyclophosphamide). Medications to manage symptoms and address underlying mechanisms Lifestyle modifications and self-care measures Physical therapy and rehabilitation Surgical intervention when indicated Regular monitoring and follow-up care Patient education and support services

\begin{itemize}[noitemsep]
  \item Medications to manage symptoms and address underlying mechanisms
  \item Lifestyle modifications and self-care measures
  \item Physical therapy and rehabilitation
  \item Surgical intervention when indicated
  \item Regular monitoring and follow-up care
\end{itemize}
\section{Living with It}
\begin{itemize}[noitemsep]
  \item Follow your treatment plan as prescribed
  \item Maintain regular follow-up with your healthcare provider
  \item Make necessary lifestyle adjustments
  \item Seek support from family, friends, or support groups
  \item Stay informed about your condition
\end{itemize}
\section{Outlook}
The outlook has improved dramatically: mortality declined from >90\% to 10--20\% with prompt PLEX. The outlook varies depending on the specific condition, severity at diagnosis, and response to treatment. Early intervention generally leads to better outcomes. Many conditions can be effectively managed with appropriate treatment. Regular follow-up care is important for monitoring and treatment adjustment.
